% !TEX program = xelatex
% AI-effects 프로젝트 — AEI 사용량 데이터를 활용한 실증연구 11편이 "사용량"을
% Claude.ai 데이터와 1P API 데이터 중 무엇으로 정의했는지(합산/단독/병렬)를
% 원문 PDF(부록 수식 포함)와 원자료 문서로 대조해 정리한 작업 노트(report/).
% 컴파일: xelatex -> biber -> xelatex -> xelatex (kotex, biblatex/biber 필요)
\documentclass[11pt]{article}

\usepackage{fontspec}
\usepackage{kotex}
\usepackage[a4paper,margin=2.2cm]{geometry}
\usepackage{longtable}
\usepackage{booktabs}
\usepackage{array}
\usepackage{amsmath}
\usepackage{amssymb}
\usepackage{xcolor}
\usepackage{hyperref}
\usepackage{enumitem}
\usepackage{titlesec}
\usepackage[backend=biber,style=authoryear,maxcitenames=2,uniquelist=false,sorting=nyt]{biblatex}
\addbibresource{../../reference/references.bib}

\hypersetup{colorlinks=true,linkcolor=blue!50!black,citecolor=blue!50!black,urlcolor=blue!50!black}
\newcolumntype{L}[1]{>{\raggedright\arraybackslash}p{#1}}
\setlength{\LTleft}{0pt}
\setlength{\LTright}{0pt}
\renewcommand{\arraystretch}{1.15}
\setlength{\tabcolsep}{4pt}

\title{AEI 실증연구는 ``사용량''을\\[1mm]
어느 플랫폼으로 정의했는가\\[2mm]
\large --- Claude.ai 데이터와 1P API 데이터를 합산했는가, 하나만 썼는가, 따로 비교했는가 ---}
\author{AI-effects 프로젝트}
\date{\today}

\begin{document}
\maketitle

\begin{abstract}
\noindent Anthropic Economic Index(AEI)의 사용량(usage) 데이터는 두 갈래로 공개된다. (1) 소비자용 대화 서비스 \textbf{Claude.ai}(Free$\cdot$Pro$\cdot$Max, V6부터 Cowork 포함)의 대화 표본과 (2) 개발자$\cdot$기업이 Anthropic API를 직접 호출하는 \textbf{1P(first-party) API} 트래픽 표본이다. 본 노트는 AEI 사용량 데이터를 실증에 활용한 연구 11편이 ``사용량''을 두 데이터 중 무엇으로 정의했는지를 원문(PDF)과 원자료 문서로 대조해 정리한다. 결론은 세 가지다. 첫째, 11편 중 5편(\cite{handa-2025-which-economic-tasks-are}; \cite{tamkin-2025-estimating-ai-productivity-gains-from}; \cite{fan-2026-what-countries-use-ai-and}; \cite{fan-2026-aggregate-gains-from-ai}; \cite{bick-2026-what-work-does-generative})은 \textbf{Claude.ai만} 사용량으로 정의했다. 둘째, AEI 정기 보고서 4편(V3--V6)은 두 데이터를 \textbf{합치지 않고 플랫폼별로 따로 측정해 비교}했으며, 그 안에서도 지리 지표$\cdot$실효 AI 커버리지$\cdot$tenure 분석처럼 Claude.ai에만 존재하는 지표가 많다. 셋째, 두 데이터를 \textbf{하나의 사용량으로 합산한 연구는 \textcite{massenkoff-2026-labor-market-impacts-of-ai}의 관측노출(observed exposure) 하나뿐}이며, 이 합산은 대칭적이지 않다---Claude.ai는 업무용 대화만, API는 전체 트래픽을 세고, 자동화 가중치에서는 API 사용 전체를 ``자동화''로 간주한다. \textcite{kim-2026-generative-ai-korean-labor-market}은 이 관측노출을 한국 직업분류에 그대로 이식했으므로 합산 정의를 간접적으로 계승한다.
\end{abstract}

\tableofcontents

% ============================================================
\section{서론: 질의와 확인 방법}
% ============================================================

질의: AEI 사용량 데이터를 활용한 각 연구가 사용량을 어떻게 정의했는지, 특히 \textbf{Claude.ai 데이터와 1P API 데이터를 합쳐서 정의했는지, 둘 중 하나만 정의했는지}를 연구별로 정리한다.

확인 방법:
\begin{enumerate}[label=(\arabic*)]
  \item 대상 선정: report 07(\texttt{07\_AEI실증연구\_노출도사용량\_계보.tex})이 확정한 AEI 직접 활용 연구 9편에, AEI 과업 점유율을 비교 벤치마크로 쓰는 \textcite{bick-2026-what-work-does-generative}와 AEI 관측노출을 한국에 적용한 \textcite{kim-2026-generative-ai-korean-labor-market}을 더해 11편.
  \item 각 연구의 \texttt{reference/notes/} 노트로 1차 분류한 뒤, 판단이 갈리는 부분은 \texttt{reference/papers/}의 원문 PDF를 텍스트 추출(pdftotext)해 ``Claude.ai'', ``1P API'', ``first-party'' 언급을 전수 검색했다.
  \item \textcite{massenkoff-2026-labor-market-impacts-of-ai} 부록의 사용량 정의식은 PDF 안에 이미지로 들어 있어 텍스트 추출이 되지 않으므로, 해당 페이지(부록 p.2--3)를 이미지로 렌더링해 수식을 직접 확인했다.
  \item 공개 데이터의 플랫폼 구분은 \texttt{data/raw/usage/anthropic\_economic\_index/}에 내려받은 각 릴리스의 \texttt{README.md}$\cdot$\texttt{data\_documentation.md}로 확인했다.
\end{enumerate}

% ============================================================
\section{두 데이터의 정의와 공개 데이터 구조}
\label{sec:platform}
% ============================================================

\textbf{Claude.ai}: 소비자용 채팅 서비스 대화. V6 문서 기준 ``Claude chat and Cowork (Free, Pro, and Max plans)''. IP 주소 기반으로 국가$\cdot$미국 주$\cdot$하위지역이 식별되어 \textbf{지리 분해가 가능}하며, 대화 목적(업무/코스워크/개인) 분류가 붙는다.

\textbf{1P API}: Anthropic API를 직접 호출하는 개발자$\cdot$기업 트래픽(Amazon Bedrock$\cdot$Google Vertex 등 제3자 경유 트래픽은 제외, V6부터 Claude Code 제외). 각 레코드는 단일 입력--출력 쌍이며, \textbf{지리 정보가 없어 전세계(GLOBAL) 집계만 공개}된다.

\begin{longtable}{L{4.2cm}L{2.6cm}L{2.6cm}L{6.4cm}}
\caption{릴리스별 플랫폼 공개 현황(원자료 문서 기준)}\label{tab:release}\\
\toprule
\textbf{릴리스} & \textbf{Claude.ai} & \textbf{1P API} & \textbf{비고}\\
\midrule
\endfirsthead
\toprule
\textbf{릴리스} & \textbf{Claude.ai} & \textbf{1P API} & \textbf{비고}\\
\midrule
\endhead
\bottomrule
\endfoot
\texttt{release\_2025\_02\_10} (R1) & ○ & $\times$ & 과업별 사용 비중(\texttt{onet\_task\_mappings.csv})\\
\texttt{release\_2025\_03\_27} (V2) & ○ & $\times$ & \texttt{task\_pct\_v1/v2.csv}\\
\texttt{release\_2025\_09\_15} (V3) & ○ (국가$\cdot$미국 주) & ○ (GLOBAL만) & API 최초 공개, 파일 분리\\
\texttt{release\_2026\_01\_15} (V4) & ○ & ○ & 플랫폼별 raw 파일 분리\\
\texttt{release\_2026\_03\_24} (V5) & ○ & ○ & 플랫폼별 raw 파일 분리\\
\texttt{release\_2026\_06\_26} (V6) & ○ (국가$\cdot$하위지역) & ○ (GLOBAL만) & \texttt{aei\_claude\_ai\_*.csv} / \texttt{aei\_1p\_api\_*.csv}\\
\texttt{labor\_market\_impacts} & 합산 & 합산 & \texttt{job\_exposure.csv}, \texttt{task\_penetration.csv}\\
\end{longtable}

AEI 정기 릴리스는 한 번도 두 플랫폼을 합친 사용량 파일을 공개한 적이 없다. 두 플랫폼이 합산된 공개 산출물은 \textcite{massenkoff-2026-labor-market-impacts-of-ai}의 결과물인 \texttt{labor\_market\_impacts/} 폴더뿐이다.

% ============================================================
\section{연구별 사용량 정의}
\label{sec:studies}
% ============================================================

\begin{longtable}{L{3.4cm}L{4.2cm}L{2.2cm}L{2.6cm}L{3.4cm}}
\caption{연구별 사용량 정의(Claude.ai / 1P API)}\label{tab:studies}\\
\toprule
\textbf{연구} & \textbf{사용 데이터} & \textbf{Claude.ai} & \textbf{1P API} & \textbf{사용량 정의 유형}\\
\midrule
\endfirsthead
\toprule
\textbf{연구} & \textbf{사용 데이터} & \textbf{Claude.ai} & \textbf{1P API} & \textbf{사용량 정의 유형}\\
\midrule
\endhead
\bottomrule
\endfoot
\textcite{handa-2025-which-economic-tasks-are} [R1] & 2024.12--2025.1, Claude.ai Free$\cdot$Pro 100만 건(전체 400만+) & 사용 & 미사용(엔터프라이즈$\cdot$API 제외 명시) & A. Claude.ai 단독\\
\textcite{tamkin-2025-estimating-ai-productivity-gains-from} & Claude.ai Free$\cdot$Pro$\cdot$Max 10만 건 & 사용 & 미사용 & A. Claude.ai 단독\\
\textcite{appel-2025-uneven-geographic-and-enterprise-ai} [V3] & 2025.8, Claude.ai 100만 + 1P API 100만 & 사용(과업$\cdot$지리) & 사용(3장 기업 사용) & B. 플랫폼별 병렬 비교\\
\textcite{appel-2026-economic-primitives} [V4] & 2025.11, Claude.ai 100만 + 1P API 100만 & 사용 & 사용 & B. 플랫폼별 병렬 비교\\
\textcite{massenkoff-2026-labor-market-impacts-of-ai} & V3$\cdot$V4(2025.8$\cdot$11), Claude.ai 200만 + 1P API 200만 & 업무용 대화만 & 전체 트래픽 & \textbf{C. 합산}\\
\textcite{massenkoff-2026-learning-curves} [V5] & 2026.2, Claude.ai 100만 + 1P API 100만 & 사용 & 사용 & B. 플랫폼별 병렬 비교\\
\textcite{massenkoff-2026-cadences} [V6] & 2026.4--6, Claude 대화(chat$\cdot$Cowork), Claude Code, 1P API & 사용 & 사용 & B. 병렬 비교(+Claude Code). 3장은 C형 관측노출 차용\\
\textcite{fan-2026-what-countries-use-ai-and} & V5 릴리스, Claude.ai 약 80만 건(국가 식별분) & 사용 & 미사용(국가별 미공개) & A. Claude.ai 단독\\
\textcite{fan-2026-aggregate-gains-from-ai} & R1--R5, Claude.ai & 사용 & 각주의 보조 추정에만 사용 & A. Claude.ai 단독\\
\textcite{bick-2026-what-work-does-generative} & V2 릴리스(2025.3.27) 과업 점유율 & 사용(비교용) & 미사용 & A. Claude.ai 단독\\
\textcite{kim-2026-generative-ai-korean-labor-market} & \textcite{massenkoff-2026-labor-market-impacts-of-ai} 관측노출 & 간접 & 간접 & C. 합산(계승)\\
\end{longtable}

\subsection{유형 A: Claude.ai만 사용량으로 정의}

\begin{itemize}
  \item \textbf{\textcite{handa-2025-which-economic-tasks-are}}: AEI 최초 보고서. 사용량 = Claude.ai Free$\cdot$Pro 대화의 O*NET 과업별 점유율. 자료 설명에서 ``엔터프라이즈/API 고객 데이터 제외''를 명시하고, 한계 절에서 API 데이터와는 사용자층$\cdot$제품 기능이 달라 대표성이 제한된다고 적는다. ``직업의 36\%가 과업의 25\% 이상에서 AI를 사용'' 같은 커버리지 수치도 Claude.ai만으로 계산된 것이다.
  \item \textbf{\textcite{tamkin-2025-estimating-ai-productivity-gains-from}}: Claude.ai(Free$\cdot$Pro$\cdot$Max) 대화 10만 건으로 과업별 ``AI 없이/AI와 함께'' 소요시간을 추정하고 이를 Hulten 정리로 집계(연 1.8\%p 생산성 효과). 원문은 사용 표본이 Claude.ai 대화로만 구성되어 AI 활용 전체를 대표하지 못한다고 명시한다.
  \item \textbf{\textcite{fan-2026-what-countries-use-ai-and}}: V5 릴리스(2026.2)로 국가별 사용 강도(인구 백만 명당 대화 수)와 확산폭(요청군집 HHI)을 측정. 1P API는 국가별로 공개되지 않아 \textbf{구조적으로 쓸 수 없었다}고 밝히고, API 스트림의 국가별 공개를 향후 과제로 제시한다. API는 소프트웨어 집중도(컴퓨터$\cdot$수학 51.6\%)를 보여주는 전세계 수치로만 각주에서 인용된다.
  \item \textbf{\textcite{fan-2026-aggregate-gains-from-ai}}: R1--R5 다섯 웨이브의 Claude.ai 데이터로 노동비용등가(LCE, 연 약 2.7조 달러)와 AI 집중지수(ACI)를 계산. R5 API 표본을 쓰면 API만으로 연 약 1.14조 달러의 잠재 LCE가 나온다는 보조 추정을 각주에 제시하지만, \textbf{모든 공식 추정치에서 API는 제외}된다. 저자는 API 트래픽이 컴퓨터$\cdot$수학직에 치우쳐 있어 Claude.ai만 쓴 ACI가 실제 집중도의 하한일 수 있다고 적는다.
  \item \textbf{\textcite{bick-2026-what-work-does-generative}}: 자체 서베이(RPS)가 주 자료이고, AEI는 비교 벤치마크다. 비교에 쓴 것은 ``2025년 3월 27일 AEI 릴리스''(V2)의 과업 점유율로, 이 릴리스에는 Claude.ai 데이터만 있다(표~\ref{tab:release}). 원문 각주는 이후 OpenAI$\cdot$Anthropic이 기업 사용 데이터를 공개했다고만 언급한다.
\end{itemize}

\subsection{유형 B: 두 데이터를 합치지 않고 플랫폼별로 병렬 비교}

AEI 정기 보고서 V3--V6은 매번 Claude.ai와 1P API를 각각 100만 건씩 표집하지만, \textbf{두 표본을 하나의 사용량으로 합치지 않는다}. 모든 그림$\cdot$표가 ``Claude.ai vs 1P API'' 두 패널로 나뉘고, 결론도 두 플랫폼의 차이(예: API는 코딩$\cdot$자동화 비중이 훨씬 높음)를 보여주는 방식이다. 다만 보고서 안의 개별 지표 중에는 \textbf{Claude.ai에만 정의된 것이 많다}는 점이 중요하다.

\begin{itemize}
  \item \textbf{\textcite{appel-2025-uneven-geographic-and-enterprise-ai} [V3]}: 1--2장(과업 구성 변화, 국가$\cdot$미국 주별 Anthropic AI Usage Index(AUI))은 Claude.ai만, 3장(기업 사용, API 비용 지수, 토큰 탄력성)은 1P API만 쓴다. 과업 집중도(지니계수 0.84 vs 0.86)나 자동화 비중(Claude.ai 약 50\% vs API 77\%)처럼 두 플랫폼을 나란히 놓고 비교할 뿐이다. 원문은 ``지리적 사용 패턴은 현재 Claude.ai 트래픽에만 존재한다''고 명시한다.
  \item \textbf{\textcite{appel-2026-economic-primitives} [V4]}: 경제 원시지표(과업 복잡성, 교육 수준, 용도, 자율성, 성공률)를 두 플랫폼에 각각 적용해 비교한다(예: 업무 비중 API 74\% vs Claude.ai 46\%). 생산성 추정도 플랫폼별로 따로 하는데, 성공률을 반영하지 않으면 둘 다 연 1.8\%p이고 반영하면 Claude.ai 1.2\%p, API 1.0\%p다. 반면 \textbf{실효 AI 커버리지(effective AI coverage)는 Claude.ai만으로 정의}된다(그림 4.4 설명: ``based on Claude.ai data. Task coverage is the share of tasks that appear in Claude.ai usage''). 지리 분석도 Claude.ai만 쓴다.
  \item \textbf{\textcite{massenkoff-2026-learning-curves} [V5]}: 상위 10대 과업 점유율(Claude.ai 24\%$\to$19\%, API 28\%$\to$33\%), 평균 과업 가치, 협업 모드를 플랫폼별로 비교한다. 모델 선택(Opus 비중)--임금 관계도 Claude.ai(+1.48\%p)와 API(+2.79\%p)를 따로 추정한다. tenure(가입 후 경과 기간) 분석은 가입 계정 단위 개념이라 Claude.ai 이용자 대상이다. 원문은 ``signed up for Claude''라고만 쓰며 플랫폼을 명시하지는 않는다.
  \item \textbf{\textcite{massenkoff-2026-cadences} [V6]}: 분석 단위를 ``Claude 대화(Claude.ai$\cdot$Cowork)'', ``Claude Code'', ``1P API'' 셋으로 나눠 따로 보고한다(요일$\cdot$시간대 리듬, 산출물 분류). Claude Code가 제3의 축으로 추가되었는데, 공개 데이터(V6)에서는 1P API에서도 Claude Code가 제외되어 있다. 3장 설문 분석은 선행연구의 관측노출을 비교 변수로 쓰므로 그 부분만 유형 C 정의를 차용한다.
\end{itemize}

\subsection{유형 C: 두 데이터를 합산해 하나의 사용량으로 정의}
\label{sec:typeC}

\textbf{\textcite{massenkoff-2026-labor-market-impacts-of-ai}}는 두 플랫폼을 하나의 사용량으로 합친 유일한 연구다. 부록 p.2--3의 정의식은 다음과 같다.

\begin{align}
\text{WorkUsage}_t &= \text{ClaudeWorkUsage}_t + \text{APIUsage}_t \label{eq:workusage}\\
\tilde{r}_t &= \mathbb{1}\{\text{WorkUsage}_t \ge 100\} \times \mathbb{1}\{\beta_t \ge 0.5\} \times \alpha_t \label{eq:rt}\\
\alpha_t &= \frac{1}{2} + \frac{1}{2}\times\frac{\text{ClaudeWorkUsage}_t \times \text{AutoShare}_t + \text{APIUsage}_t}{\text{ClaudeWorkUsage}_t + \text{APIUsage}_t} \label{eq:alpha}
\end{align}

\begin{itemize}
  \item $\text{ClaudeWorkUsage}_t$: Claude.ai 대화 중 \textcite{appel-2026-economic-primitives}의 용도 분류상 \textbf{업무 관련(work-related)}으로 분류된 것만 센다. 코스워크$\cdot$개인용은 제외한다.
  \item $\text{APIUsage}_t$: 1P API 트래픽은 업무 여부를 구분하지 않고 \textbf{전부} 센다. API 호출은 대개 생산 워크플로에 통합된 사용이기 때문이라고 설명한다.
  \item 커버리지 관문: $\text{WorkUsage}_t \ge 100$(직전 두 보고서 Claude.ai 200만 + API 200만 건 기준으로 전체 트래픽의 0.0025\%)을 넘어야 과업이 ``커버됨''으로 인정된다.
  \item $\text{AutoShare}_t$는 \textbf{Claude.ai 사용 중} 자동화형 비중이다. 식~\eqref{eq:alpha}의 분자에 $\text{APIUsage}_t$가 그대로 더해지므로 \textbf{API 사용은 전부 자동화형으로 간주}된다. 원문 예시로, 의료정보 코딩 과업은 사용의 90\% 이상이 1P API에서 나와 $\alpha$가 0.96이 된다.
\end{itemize}

이 합산은 단순한 ``두 표본의 합''이 아니며, 세 가지 비대칭이 있다.
\begin{enumerate}[label=(\arabic*)]
  \item \textbf{범위의 비대칭}: Claude.ai는 업무용만, API는 전량을 센다. 부록은 이를 ``API 사용에 더 큰 가중치를 준다(more weight given to API usage)''고 표현한다.
  \item \textbf{자동화 판정의 비대칭}: Claude.ai는 대화별 협업 모드 분류로 자동화 비중을 계산하지만, API는 분류 없이 전량 자동화로 처리한다.
  \item \textbf{결과적 효과}: API 트래픽이 많은 과업(코딩, 고객 응대, 데이터 입력 등)은 관문 통과와 자동화 가중치 두 경로 모두에서 노출도가 올라간다. 본문도 고객서비스 상담원이 2위(70.1\%)인 이유로 ``주요 과업이 1P API 트래픽에서 점점 더 많이 관측되기 때문''을 든다.
\end{enumerate}

참고: \texttt{reference/notes/massenkoff-2026-labor-market-impacts-of-ai.md}는 $\alpha_t$를 ``$1/2 + 1/2\times\text{AutomationShare}_t$''로만 요약해, API 사용이 자동화 비중의 분자에 들어간다는 점이 빠져 있다. 위 식~\eqref{eq:alpha}는 부록 원문 p.3에서 직접 확인한 것이다.

\textbf{\textcite{kim-2026-generative-ai-korean-labor-market}}은 \textcite{massenkoff-2026-labor-market-impacts-of-ai}의 직업별 관측노출을 SOC 2018 $\to$ ISCO $\to$ KSCO 8차 세분류로 연계해 지역별 고용조사와 결합한다. 사용량을 새로 정의하지 않고 관측노출 값을 그대로 이식하므로 합산 정의(유형 C)를 간접적으로 계승한다. 원문은 이 지표를 ``Claude 실제 사용 기록''으로만 서술하며, 1P API 트래픽이 합산되어 있다는 점이나 API가 전량 자동화로 가중된다는 점은 언급하지 않는다.

% ============================================================
\section{종합: 무엇을 사용량으로 삼았는가}
% ============================================================

\begin{longtable}{L{2.4cm}L{4.4cm}L{4.6cm}L{4.4cm}}
\caption{유형별 요약}\label{tab:summary}\\
\toprule
\textbf{유형} & \textbf{정의} & \textbf{해당 연구} & \textbf{이유$\cdot$특징}\\
\midrule
\endfirsthead
\toprule
\textbf{유형} & \textbf{정의} & \textbf{해당 연구} & \textbf{이유$\cdot$특징}\\
\midrule
\endhead
\bottomrule
\endfoot
A. Claude.ai 단독 & Claude.ai 대화의 과업$\cdot$직업$\cdot$국가별 점유율 또는 건수 & \textcite{handa-2025-which-economic-tasks-are}; \textcite{tamkin-2025-estimating-ai-productivity-gains-from}; \textcite{fan-2026-what-countries-use-ai-and}; \textcite{fan-2026-aggregate-gains-from-ai}; \textcite{bick-2026-what-work-does-generative} & 초기 릴리스에 API가 없었거나(R1$\cdot$V2), 국가 단위 분석에 API를 쓸 수 없음(GLOBAL만 공개)\\
B. 병렬 비교 & 두 플랫폼을 각각 측정하고 합치지 않음 & \textcite{appel-2025-uneven-geographic-and-enterprise-ai}; \textcite{appel-2026-economic-primitives}; \textcite{massenkoff-2026-learning-curves}; \textcite{massenkoff-2026-cadences} & 소비자 사용과 기업 배치의 차이 자체가 분석 대상. 지리$\cdot$실효 커버리지$\cdot$tenure 지표는 Claude.ai 전용\\
C. 합산 & WorkUsage = Claude.ai 업무용 + API 전체, API는 전량 자동화 가중 & \textcite{massenkoff-2026-labor-market-impacts-of-ai}; (계승) \textcite{kim-2026-generative-ai-korean-labor-market} & 노동시장 영향 측정 목적상 ``업무 관련 사용''을 넓게 포착하려는 설계\\
\end{longtable}

11편 중 두 데이터를 합산한 것은 사실상 1개 지표(관측노출)뿐이며, 나머지는 모두 Claude.ai 단독이거나 플랫폼별로 분리되어 있다. Claude.ai 단독 연구가 많은 가장 큰 이유는 데이터 구조다. 1P API는 전세계 합계만 공개되므로 \textbf{국가$\cdot$지역 단위 분석에서는 Claude.ai만이 유일한 선택지}다.

% ============================================================
\section{본 프로젝트에 대한 시사점}
% ============================================================

\begin{enumerate}[label=(\arabic*)]
  \item \textbf{한국 단위 사용량은 Claude.ai로만 정의 가능하다.} V6의 \texttt{geo\_id="KOR"} 행은 \texttt{aei\_claude\_ai\_*.csv}에만 존재하고 1P API는 GLOBAL뿐이다(report 09). 따라서 한국 사용량 지표는 유형 A가 될 수밖에 없으며, 한국 기업의 API 배치는 포착하지 못한다는 한계를 명시해야 한다.
  \item \textbf{\texttt{job\_exposure.csv}를 쓰면 유형 C가 된다.} \texttt{labor\_market\_impacts/job\_exposure.csv}와 \texttt{task\_penetration.csv}는 Claude.ai 업무용과 API 전체를 합산하고 API를 전량 자동화로 가중한 값이다. \textcite{kim-2026-generative-ai-korean-labor-market}처럼 이를 한국에 적용하면, 미국 중심의 전세계 API 배치 패턴이 한국 직업 노출도에 섞여 들어간다.
  \item \textbf{두 유형을 섞어 비교할 때 주의해야 한다.} 한국 Claude.ai 사용량(유형 A)과 관측노출(유형 C)을 한 모형에서 함께 쓰거나 비교할 경우, 두 지표의 차이 중 일부는 ``한국 vs 미국''이 아니라 ``API 포함 여부''에서 온다.
  \item \textbf{재현 가능한 대안.} V6 GLOBAL에서 Claude.ai와 1P API의 과업별 사용량을 각각 가져오면, \textcite{massenkoff-2026-labor-market-impacts-of-ai}의 합산식을 (i) Claude.ai만, (ii) 합산 두 버전으로 다시 계산해 API 포함 여부의 민감도를 점검할 수 있다. 다만 공개된 \texttt{task\_penetration.csv}에는 과업 ID$\cdot$직업 코드가 없고, 가중치 $w_t$(과업 시간 비중)도 공개되지 않아 원 지표를 정확히 재현하려면 별도 작업이 필요하다.
\end{enumerate}

% ============================================================
\section{보론: GLOBAL 단위에서 두 플랫폼 사용량을 합산할 수 있는가}
\label{sec:globalsum}
% ============================================================

질의: \texttt{geo\_id="GLOBAL"} 수준에서 \texttt{1p\_api}의 사용량과 \texttt{claude\_ai}의 사용량을 합산할 수 있는가. 1P API는 GLOBAL만 공개되므로(\S\ref{sec:platform}) 두 플랫폼을 함께 쓸 수 있는 곳은 GLOBAL 행뿐이다. 아래는 V6(\texttt{release\_2026\_06\_26}) 원자료 두 파일을 직접 집계해 확인한 결과다.

결론부터 말하면, \textbf{기술적으로는 결합할 수 있지만 두 플랫폼의 상대 규모를 알려주는 정보가 공개되어 있지 않다. 따라서 합산 가중치를 연구자가 가정으로 정해야 한다.} 합산 사용량은 관측된 양이 아니라 가정에 의존하는 구성 지표다.

\subsection{결합이 가능한 근거}

\begin{itemize}
  \item 두 파일의 GLOBAL 행은 분류 체계(\texttt{onet} 과업$\to$DWA$\to$IWA$\to$GWA, \texttt{soc\_occupation}, \texttt{request})와 지표(\texttt{pct}, \texttt{use\_case\_work\_pct}, \texttt{collaboration\_bucket\_automation\_pct} 등)가 같다. 따라서 \texttt{date\_start}$\times$\texttt{node\_external\_id}로 1:1 결합할 수 있다.
  \item 두 파일 모두 2026년 4월$\cdot$5월 두 달을 담고 있어 기간이 일치한다.
\end{itemize}

\subsection{단순 합산이 안 되는 이유}

\begin{enumerate}[label=(\arabic*)]
  \item \textbf{분모가 다르다.} V6의 \texttt{pct}는 각 플랫폼 안에서 정규화된 점유율이다. GLOBAL \texttt{overall}의 \texttt{usage\_pct}는 두 파일 모두 100.0이고, 대화 건수$\cdot$토큰량 같은 규모 지표는 없다. ``Claude.ai 트래픽이 1P API의 몇 배인가''를 데이터로 알 수 없으므로, 과업 $t$의 합산 점유율은 다음과 같이 가중치 $w$를 외생적으로 정해야만 정의된다.
  \begin{equation}
  \text{pct}^{\text{pooled}}_t(w) = w\cdot \text{pct}^{\text{claude\_ai}}_t + (1-w)\cdot \text{pct}^{\text{1p\_api}}_t,\qquad 0\le w\le 1 \label{eq:pooled}
  \end{equation}
  \item \textbf{이전 릴리스의 count도 해결책이 아니다.} V3--V5의 raw 파일(\texttt{aei\_raw\_*.csv})에는 \texttt{onet\_task\_count} 같은 건수 변수가 있다. 그러나 이는 플랫폼마다 약 100만 건씩 따로 뽑은 표본의 건수이며 실제 사용 규모를 반영하지 않는다. \textcite{massenkoff-2026-labor-market-impacts-of-ai}의 $\text{WorkUsage}_t = \text{ClaudeWorkUsage}_t + \text{APIUsage}_t$(식~\eqref{eq:workusage})도 두 표본의 크기가 같으므로(각 200만 건), 사실상 식~\eqref{eq:pooled}에서 \textbf{표본 동일 가중}($w\approx 0.5$, 업무용 필터 적용 전 기준)을 암묵적으로 가정한 것이다.
  \item \textbf{과업 단위 \texttt{pct}는 합이 100이 아니고, 공개 과업 집합도 다르다.} 표~\ref{tab:globalsum}처럼 GLOBAL 과업(레벨 0) \texttt{pct}의 합은 100에 못 미친다. 과업에 배정되지 않았거나 공개 임계치에 미달한 부분이 있다는 뜻이다. 두 플랫폼에서 공개된 과업 집합도 일부만 겹친다. 문서는 ``누락된 행은 값이 0이라는 뜻이 아니다''라고 명시하므로, 한쪽에만 있는 과업의 다른 쪽 값을 0으로 채우는 것도 가정이다. 공개 과업끼리 다시 정규화할지도 함께 정해야 한다.
  \item \textbf{두 플랫폼의 성격이 크게 다르다.} 업무용 비중과 자동화 비중이 두 배 가까이 차이 난다(표~\ref{tab:globalsum}). 그래서 합산 결과의 과업 구성이 $w$에 민감하게 달라진다.
\end{enumerate}

\begin{longtable}{L{6.2cm}L{2.2cm}L{2.2cm}L{2.2cm}L{2.2cm}}
\caption{V6 GLOBAL 행의 플랫폼별 비교(원자료 직접 집계)}\label{tab:globalsum}\\
\toprule
\textbf{항목} & \textbf{claude\_ai 2026.4} & \textbf{1p\_api 2026.4} & \textbf{claude\_ai 2026.5} & \textbf{1p\_api 2026.5}\\
\midrule
\endfirsthead
\toprule
\textbf{항목} & \textbf{claude\_ai 2026.4} & \textbf{1p\_api 2026.4} & \textbf{claude\_ai 2026.5} & \textbf{1p\_api 2026.5}\\
\midrule
\endhead
\bottomrule
\endfoot
\texttt{overall} \texttt{usage\_pct} & 100.0 & 100.0 & 100.0 & 100.0\\
\texttt{overall} \texttt{use\_case\_work\_pct} (\%) & 45.53 & 82.19 & 43.36 & 84.03\\
\texttt{overall} \texttt{collaboration\_bucket\_automation\_pct} (\%) & 48.98 & 93.66 & 48.62 & 94.22\\
과업(레벨 0) \texttt{pct} 합계 (\%) & 88.40 & 81.13 & 94.32 & 85.00\\
공개 과업 수 & 2,410 & 1,992 & 2,713 & 2,295\\
두 플랫폼 공통 과업 수 & \multicolumn{2}{l}{1,577 (합집합 2,825)} & \multicolumn{2}{l}{1,815 (합집합 3,193)}\\
\end{longtable}

\subsection{권고: 합산은 민감도 분석용으로}

\begin{enumerate}[label=(\arabic*)]
  \item \textbf{기본 사양은 플랫폼별 분리(유형 B)로 둔다.} 합산 지표는 본 분석이 아니라 민감도 분석용으로 만든다.
  \item 합산할 때는 적어도 다음 세 버전을 함께 제시한다. (i) Claude.ai 단독($w=1$), (ii) 표본 동일 가중($w=0.5$, \textcite{massenkoff-2026-labor-market-impacts-of-ai}와 같은 암묵적 가정), (iii) \textcite{massenkoff-2026-labor-market-impacts-of-ai}식 업무용 필터 버전. (iii)은 Claude.ai 쪽을 $\text{pct}^{\text{claude\_ai}}_t\times\text{use\_case\_work\_pct}_t/100$로 줄이고 1P API는 전량 반영한다. 여기에 $w$를 0에서 1까지 바꿔 가며 주요 결과가 얼마나 달라지는지 보고한다.
  \item 한쪽 플랫폼에만 공개된 과업을 어떻게 처리했는지(0 대체 여부, 재정규화 여부)를 do 파일 주석과 \texttt{log.md}에 명시한다.
  \item 어떤 버전이든 GLOBAL 합산값은 \textbf{전세계 직업$\cdot$과업 구성 지표}일 뿐이며, 한국 사용 강도(\texttt{geo\_id="KOR"}, Claude.ai 전용; report 09)와는 성격이 다른 지표임을 명시한다.
\end{enumerate}

% ============================================================
\printbibliography[title={참고문헌}]
% ============================================================

\end{document}
